\section{Multibasin benchmark inventory and system definitions}
\label{app:fixed17_inventory}

\Cref{tab:fixed17_inventory} lists the \(15\) two-dimensional multibasin systems used in this paper. 
Every row contributes to the retained-system aggregate in
\cref{tab:fixed17_alignment}, \cref{tab:self_routing}, and
\cref{fig:fixed17_horizon_curves}. Four systems are additionally displayed as
the main benchmark visual panels in \cref{fig:benchmarks}. The support-refresh
subtable uses \(12\) retained systems for which the controlled-transfer
construction produced eligible post-entry comparisons. 

Basin labels are used
only for benchmark annotation, controlled-transfer construction, and evaluation.
Basin counts are benchmark metadata; for the
structured-transition ablation rows, they are also used only to set a fixed
architecture group count, as described in \cref{app:experimental_details}, and
not for support extraction, support-family construction, routing, or rollout
evaluation.

The 15 multibasin benchmark systems were generated procedurally with Claude Opus 4.5. 

\begin{table}[tbp]
\centering
\caption{The \(15\) multibasin benchmark systems.}
\label{tab:fixed17_inventory}
\resizebox{\textwidth}{!}{%
\begin{tabular}{@{}l l l c@{}}
\toprule
Display name & System key & Generator family & Basins \\
\midrule
Local-linear gates
  & \texttt{gated\_local\_linear}
  & piecewise local-linear gates & 3 \\
Transfer-gated local-linear
  & \texttt{gated\_transfer\_linear}
  & piecewise transfer gates & 3 \\
Arrested spiral
  & \texttt{arrested\_spiral}
  & spiral capture with Gaussian traps & 5 \\
Asymmetric three-well
  & \texttt{cal\_asymmetric\_3}
  & Gaussian wells with independent rotation & 3 \\
High-cross three-well
  & \texttt{cal\_high\_cross\_3}
  & Gaussian wells with stronger rotation & 3 \\
Hexagonal six-well
  & \texttt{cal\_hexagon\_6}
  & Gaussian wells with independent rotation & 6 \\
Octagonal eight-well
  & \texttt{cal\_octagon\_8}
  & Gaussian wells with independent rotation & 8 \\
Pentagonal five-well
  & \texttt{cal\_pentagon\_5}
  & Gaussian wells with independent rotation & 5 \\
Square four-well
  & \texttt{cal\_square\_4}
  & Gaussian wells with independent rotation & 4 \\
Triple-well Duffing
  & \texttt{duffing\_triple\_well}
  & triple-well Duffing oscillator & 3 \\
SNIC multi-attractor
  & \texttt{snic\_multi}
  & polar SNIC-like phase dynamics & 3 \\
Transition-routes four-well
  & \texttt{transition\_routes\_4}
  & route-augmented Gaussian wells & 4 \\
Depth-gradient four-well
  & \texttt{var\_depth\_gradient\_4}
  & Gaussian wells with depth gradient & 4 \\
Diamond four-well
  & \texttt{var\_diamond\_4}
  & Gaussian wells on a diamond layout & 4 \\
L-shaped five-well
  & \texttt{var\_l\_shape\_5}
  & Gaussian wells on an L-shaped layout & 5 \\
\bottomrule
\end{tabular}
}
\end{table}

\paragraph{Stored-step convention.}
All retained systems are implemented as continuous-time vector fields
\(\dot x=f_s(x)\) and integrated with fourth-order Runge--Kutta to produce the
stored observation map \(F_{\Delta t}\) used in \cref{eq:sampled_map}. The
training windows contain only stored states \(x_k\), not \(f_s\), integration
substeps, or basin assignments.

\paragraph{Gaussian-well family.}
Most retained catalog systems use a common two-dimensional Gaussian potential
with an independent rotational drift. For \(x=(x_1,x_2)\in\R^2\), let
\(R x=(x_2,-x_1)\) and
\[
  V_s(x)
  =
  \sum_{i=1}^{M_s}
  -a_{s,i}
  \exp\!\left(-\frac{\norm{x-c_{s,i}}_2^2}{2\sigma_{s,i}^2}\right)
  +\gamma_s(x_1^4+x_2^4).
\]
The implemented vector field is
\begin{equation}
  \dot x = -\nabla V_s(x)+\omega_s R x .
  \label{eq:app_gaussian_well}
\end{equation}
Define \(p_i^{(M)}(r,\varphi)
=r(\cos(2\pi i/M+\varphi),\sin(2\pi i/M+\varphi))\), \(i=0,\ldots,M-1\).
The retained Gaussian-well systems instantiate \cref{eq:app_gaussian_well} with
the parameters in \cref{tab:app_gaussian_params}.

\begin{table}[tbp]
\centering
\caption{Parameters for retained Gaussian-well systems. Repeated pairs in the
\((a_i,\sigma_i)\) column apply to every listed center.}
\label{tab:app_gaussian_params}
\resizebox{\textwidth}{!}{%
\begin{tabular}{@{}l l l l@{}}
\toprule
System & \(\{c_i\}\) & \((a_i,\sigma_i)\) & \((\omega,\gamma)\) \\
\midrule
High-cross three-well
  & \(\{p_i^{(3)}(1.8,\pi/2)\}\)
  & \((3.0,0.5)\) & \((2.0,0.03)\) \\
Hexagonal six-well
  & \(\{p_i^{(6)}(1.7,0)\}\)
  & \((2.0,0.5)\) & \((1.0,0.03)\) \\
Octagonal eight-well
  & \(\{p_i^{(8)}(2.2,0)\}\)
  & \((3.0,0.5)\) & \((0.90,0.02)\) \\
Pentagonal five-well
  & \(\{p_i^{(5)}(1.8,\pi/2)\}\)
  & \((2.0,0.5)\) & \((1.1,0.03)\) \\
Square four-well
  & \(\{p_i^{(4)}(1.8,\pi/4)\}\)
  & \((3.0,0.5)\) & \((1.0,0.03)\) \\
Diamond four-well
  & \(\{(0,2.2),(2.2,0),(0,-2.2),(-2.2,0)\}\)
  & \((2.5,0.5)\) & \((1.0,0.02)\) \\
L-shaped five-well
  & \(\{(-1.5,1.5),(-1.5,0),(-1.5,-1.5),(0,-1.5),(1.5,-1.5)\}\)
  & \((2.5,0.5)\) & \((1.0,0.03)\) \\
Asymmetric three-well
  & \(\{(0,1.8),(-1.5,-0.9),(1.5,-0.9)\}\)
  & \(\{(2.5,0.55),(1.5,0.4),(2.0,0.5)\}\) & \((1.0,0.03)\) \\
Depth-gradient four-well
  & \(\{(-1.3,1.3),(1.3,1.3),(1.3,-1.3),(-1.3,-1.3)\}\)
  & \(\{(2.2,0.55),(2.5,0.5),(3.0,0.5),(3.5,0.5)\}\) & \((1.3,0.03)\) \\
\bottomrule
\end{tabular}
}
\end{table}

The transition-routes system uses the same Gaussian-well form with centers
\[
  \{c_i\}
  =
  \{(-1.8,1.8),(1.8,1.8),(-1.8,-1.8),(1.8,-1.8)\},
\]
\(a_i=3.0\), \(\sigma_i=0.6\), \(\omega=1.0\), and \(\gamma=0.03\), and adds a
corridor-dependent rotation term:
\begin{equation}
  \dot x
  =
  -\nabla V_s(x)
  +
  \bigl(1.0+0.3\,C_{\rm route}(x_2)\bigr) R x,
  \qquad
  C_{\rm route}(x_2)
  =
  e^{-(x_2-1.8)^2/0.3}+e^{-(x_2+1.8)^2/0.3}.
  \label{eq:app_transition_routes}
\end{equation}

\paragraph{Native gated local-linear systems.}
Both native gated systems place three centers on a circular layout, but with
different radii and region definitions. Let
\(Q(\alpha)=\begin{psmallmatrix}\cos\alpha&-\sin\alpha\\
\sin\alpha&\cos\alpha\end{psmallmatrix}\) and
\(\alpha_b=2\pi b/3\), \(b\in\{0,1,2\}\).

For \texttt{gated\_local\_linear}, \(c_b=1.75(\cos\alpha_b,\sin\alpha_b)\).
Inside the radius-\(1.05\) core of basin \(b\),
\[
  \dot x=A_b(x-c_b),\qquad
  A_b=Q(\alpha_b)\widetilde A_b Q(\alpha_b)^\top,
\]
with
\[
\widetilde A_0=\begin{psmallmatrix}-0.9&-1.2\\1.2&-0.9\end{psmallmatrix},
\quad
\widetilde A_1=\begin{psmallmatrix}-1.35&0.2\\-0.3&-0.7\end{psmallmatrix},
\quad
\widetilde A_2=\begin{psmallmatrix}-0.7&-0.1\\0.5&-1.2\end{psmallmatrix}.
\]
Outside the cores, the active sector is the nearest angular sector \(s(x)\) and
the shared gate matrix
\[
  G=\begin{psmallmatrix}-1.35&-0.9\\0.9&-1.35\end{psmallmatrix}
\]
drives the state as \(\dot x=G(x-c_{s(x)})\).

For \texttt{gated\_transfer\_linear}, \(c_b=1.85(\cos\alpha_b,\sin\alpha_b)\).
The implemented radii and widths are
\[
\begin{aligned}
  \rho_{\rm core}&=0.30, &
  \rho_{\rm source}&=0.80, &
  \rho_{\rm exit}&=0.60, &
  \beta_{\rm exit}&=0.72,\\
  w_{\rm chan}&=0.22, &
  \delta_{\rm lane}&=0.28, &
  \rho_{\rm hand}&=0.45 .
\end{aligned}
\]
The core matrices use the same rotated form with templates
\[
\begin{psmallmatrix}-1.0&-1.1\\1.1&-1.0\end{psmallmatrix},\quad
\begin{psmallmatrix}-1.4&0.2\\-0.2&-0.7\end{psmallmatrix},\quad
\begin{psmallmatrix}-0.8&-0.3\\0.5&-1.3\end{psmallmatrix}.
\]
For each ordered pair \(p=(s,t)\), \(s\ne t\), define
\[
  d_{s\to t}=\frac{c_t-c_s}{\norm{c_t-c_s}_2},\qquad
  n_{s\to t}=(-d_{s\to t,2},d_{s\to t,1}),\qquad
  o_t=\frac{c_t}{\norm{c_t}_2},
\]
and \(\chi_{s,t}=1\) when \(\sin(\alpha_s-\alpha_t)\ge 0\), otherwise
\(\chi_{s,t}=-1\). The implemented channel entry and handoff points are
\[
\begin{aligned}
  e_p &= c_s+\rho_{\rm source}d_{s\to t}
        +\chi_{s,t}\delta_{\rm lane}n_{s\to t},\\
  h_p &= c_t+\rho_{\rm hand}o_t
        +\chi_{s,t}\delta_{\rm lane}(-o_{t,2},o_{t,1}),
\end{aligned}
\]
with channel direction \(d_p=(h_p-e_p)/\norm{h_p-e_p}_2\), channel normal
\(n_p=(-d_{p,2},d_{p,1})\), and channel length
\(\ell_p=(h_p-e_p)\cdot d_p\).

The vector field is piecewise analytic. Core regions
\(\norm{x-c_b}_2\le\rho_{\rm core}\) use \(A_b(x-c_b)\). Source annuli use
\(0.65A_b(x-c_b)\), and other non-channel background regions use
\(0.50A_b(x-c_b)\) for the nearest center \(c_b\). Exit wedges for
\(p=(s,t)\) are source-neighborhood states outside the core with
\(\norm{x-c_s}_2\ge\rho_{\rm exit}\) and
\((x-c_s)\cdot d_{s\to t}\ge\norm{x-c_s}_2\cos\beta_{\rm exit}\); they use
\[
  \dot x = v_{\rm exit}d_{s\to t}
  -\lambda_{\rm exit}\bigl((x-c_s)\cdot n_{s\to t}\bigr)n_{s\to t},
  \qquad
  v_{\rm exit}=1.0,\quad \lambda_{\rm exit}=1.8;
\]
and channel rectangles satisfying
\[
  0\le (x-e_p)\cdot d_p\le \ell_p,\qquad
  |(x-e_p)\cdot n_p|\le w_{\rm chan}
\]
use
\[
  \dot x = v_{\rm chan}d_p
  -\lambda_{\rm chan}\bigl((x-e_p)\cdot n_p\bigr)n_p,
  \qquad
  v_{\rm chan}=1.55,\quad \lambda_{\rm chan}=2.8.
\]

\paragraph{Other specialized systems.}
The arrested spiral system has a background spiral flow plus four Gaussian
traps and an origin trap:
\[
  \dot x
  =
  -0.3x+2.0Rx
  -4.0\sum_{i=1}^{4}
  e^{-\norm{x-c_i}_2^2/(2\cdot0.25)}
  \frac{x-c_i}{0.25}
  -1.5e^{-\norm{x}_2^2/0.3}\frac{x}{0.15}
  -0.02(x_1^3,x_2^3),
\]
with \(c_i\in\{(1.5,0),(0,1.8),(-1,-1),(0.8,-1.5)\}\). Its fifth basin is the
origin trap, matching the benchmark manifest count.

The triple-well Duffing system uses \(x=(q,p)\) and
\[
  V(q)=q^6/6-q^4/2+a q^2,\qquad
  V'(q)=q^5-2q^3+2a q,
\]
with \(a=0.3\), damping \(\delta=0.5\), rotation strength \(\omega=1.0\), and
soft cubic confinement \(-\epsilon u^3/s^3\) with \(\epsilon=0.003\) and
\(s=4.0\):
\begin{equation}
\begin{aligned}
  \dot q &= (1+\omega)p-\epsilon q^3/s^3,\\
  \dot p &= -V'(q)-\delta p-\omega q-\epsilon p^3/s^3 .
\end{aligned}
\label{eq:app_duffing_triple}
\end{equation}

The SNIC system is implemented in polar form and then converted to Cartesian
coordinates. With \(\varepsilon_{\rm num}=10^{-8}\),
\(r^2=x_1^2+x_2^2+\varepsilon_{\rm num}\), \(r=\sqrt{r^2}\),
\(\theta=\operatorname{atan2}(x_2,x_1)\),
\(c_\theta=x_1/(r+\varepsilon_{\rm num})\), and
\(s_\theta=x_2/(r+\varepsilon_{\rm num})\),
\[
  \dot r=r(1-r^2)-0.3\cos(3\theta),\qquad
  \dot\theta=1-1.2\cos(3\theta).
\]
The Cartesian vector field is
\[
\begin{aligned}
  \dot x_1
  &= \dot r c_\theta-r\dot\theta s_\theta
     -0.01x_1(x_1^2+x_2^2)+0.5x_2,\\
  \dot x_2
  &= \dot r s_\theta+r\dot\theta c_\theta
     -0.01x_2(x_1^2+x_2^2)-0.5x_1 .
\end{aligned}
\]
These analytic generators define the benchmark trajectories. Held-out
evaluation annotations are produced by the benchmark label helpers or by
endpoint-rollout basin labeling, and the learned models observe only stored
states.
